%% kapa: exam=2 round=23 year=2012 period=1 subject=appraisal_practice subject_ko=감정평가실무 seq=2 number=2 points=30
%% source: Q-Net 감정평가사 기출문제 (cst00309) artlSeq=5056012 file=raw/1008/5056012/1교시(감정평가실무).hwp
%% conversion: kordoc 4.24.51 -> tools/gen2.py (auto); review=proofread-sonnet; flags=none
\begin{longquestion}{【문제 2】}
선임심사역인 감정평가사 甲은 KY은행 본점에서 감정평가서 심사업무를 담당하고 있다. <자료 1>의 감정평가서를 대상으로 주어진 자료를 활용하여 다음 물음에 답하시오(단, 심사대상 감정평가서 내용 중 달리 판단할 근거가 없는 경우에는 적정한 것으로 본다). (30점) \par
\medskip\noindent 물음 1) <자료 1> 감정평가서 중 부적정한 평가내용이 있다면 구체적인 사유와 보완내용을 포함하여 기술하시오. (15점) \par
\medskip\noindent 물음 2) <자료 1> 감정평가서에서 적용하지 않은 다른 방식으로 평가가격을 검토하시오. (5점) \par
\medskip\noindent 물음 3) KY은행이 평가조건을 요구(동의)하지 않았을 경우의 감정평가액을 구하고 <자료 1> 감정평가서의 평가개요 중 달라지는 항목을 기술한 후 감정평가명세표를 재작성하시오. (10점) \par
\textbf{<자료 1> 심사대상 감정평가서(부분 요약 발췌)} \par
1. 평가개요 \par
1) 본 감정평가는 K시 H구 A동에 소재하는 부동산에 대한 담보 목적의 감정평가로 감정평가 관련법규에 따라 평가하였습니다. \par
2) 본 토지상에 사용승인을 득하고, 일반건축물대장에 등재예정인 건물(명세표상 기호 ㉮ 건물)은 귀 KY은행의 요구에 따라 일반건축물대장에 등재된 것을 전제로 하여 토지와 건물을 평가하였고 가격시점은 2012년 9월 9일입니다. \par
3) 토지는 당해 토지와 유사한 이용가치를 지닌다고 인정되는 표준지 공시지가를 기준으로 공시기준일부터 가격시점까지의 지가변동률, 당해 토지의 위치ㆍ형상ㆍ환경ㆍ이용상황 기타 가격형성상의 제요인과 인근지가수준 등을 종합적으로 참작하여 평가하였습니다. \par
4) 건물은 구조, 용도, 부대설비 및 시공상태 등을 종합 참작하여 원가법으로 평가하였으며, 감가수정은 정액법을 적용하였습니다. \par
2. 평가대상 물건현황 \par
1) 인근지역 현황 \par
대상부동산은 1990년대에 민간이 조성한 소규모 협동화공장단지 내에 소재하며 동 단지는 약 100,000m\textsuperscript{2} 규모로 13개 업체가 입주가능하고 현재의 입주율은 85\% 정도입니다. 입주업체는 의료분야의 중소규모 업체가 대부분이며 전체 종업원은 1,200명 정도입니다. 이 단지는 진입로변 국도를 통해 고속도로와 연결되어 교통, 물류 등의 여건이 비교적 양호하며 용수, 전력, 인력수급 등의 여건도 양호한 편입니다. \par
2) 대상부동산 현황 \par
당해 토지는 2차선 포장도로를 통하여 단지내 도로와 연결되고 있으며, 평탄한 콘크리트지반이 조성된 사다리형태의 상업용지로서 공부상 지목은 잡종지로 경쟁가능성이 없는 독점적인 위치를 가진 적정규모의 토지입니다. 용도지역은 동 단지 전체가 준공업지역이고, 토지의 임대상황은 없으며 등기사항전부증명서상 소유자는 乙입니다. 일반건축물대장에 소유자 乙로 등재예정인 지상 건물의 현황은 아래와 같습니다. \par
{\small\centering
\begin{tabular}{|c|c|c|c|}
\hline
\makecell{구조} & \makecell{건축연면적} & \makecell{건물규모} & \makecell{용도} \\
\hline
\makecell{철근콘크리트조 슬래브지붕} & \makecell{300m\textsuperscript{2}} & \makecell{지상 2층} & \makecell{근린생활시설} \\
\hline
\end{tabular}\par}

3. 감정평가가액 산출근거 \par
1) 토지가액 산출근거 \par
\medskip\noindent (1) 비교 표준지 공시지가: 2012.01.01 공시된 인근지역내 비교가능한 표준지는 아래와 같으며 이중 제반 여건이 유사한 기호 ② 표준지를 선정하였습니다. \par
{\small\centering
\begin{tabular}{|c|c|c|c|c|c|c|c|c|}
\hline
\makecell{기호} & \makecell{소재지} & \makecell{면적\newline{}(m\textsuperscript{2})} & \makecell{지목} & \makecell{이용\newline{}상황} & \makecell{용도지역} & \makecell{도로교통} & \makecell{형상\newline{}지세} & \makecell{공시지가\newline{}(원/m\textsuperscript{2})} \\
\hline
\makecell{①} & \makecell{A동 57} & \makecell{350} & \makecell{대} & \makecell{상업용} & \makecell{일반공업지역} & \makecell{소로한면} & \makecell{사다리평지} & \makecell{450,000} \\
\hline
\makecell{②} & \makecell{A동 154} & \makecell{630} & \makecell{장} & \makecell{주상용} & \makecell{준공업지역} & \makecell{소로한면} & \makecell{세장형평지} & \makecell{420,000} \\
\hline
\makecell{③} & \makecell{A동 322} & \makecell{245} & \makecell{대} & \makecell{단독\newline{}주택} & \makecell{준공업지역} & \makecell{소로한면} & \makecell{부정형 완경사} & \makecell{400,000} \\
\hline
\end{tabular}\par}

※ 기호 ① 표준지는 도시계획도로에 5\% 저촉됨 \par
\medskip\noindent (2) 지가변동률 \par
국토해양부장관이 조사발표한 K시 H구 공업지역의 2012.01.01부터 가격시점까지의 지가변동률은 0.750\%입니다. \par
\medskip\noindent (3) 지역요인: 인근지역에 소재하여 지역요인은 동일합니다. \par
\medskip\noindent (4) 개별요인: 대상토지가 비교표준지 ②에 비해 다소 열세합니다. \par
{\small\centering
\begin{tabular}{|c|c|c|c|c|c|c|}
\hline
\makecell{가로조건} & \makecell{접근조건} & \makecell{환경조건} & \makecell{획지조건} & \makecell{행정적 조건} & \makecell{기타조건} & \makecell{격차율} \\
\hline
\makecell{1.00} & \makecell{1.00} & \makecell{1.00} & \makecell{0.99} & \makecell{1.00} & \makecell{1.00} & \makecell{0.99} \\
\hline
\end{tabular}\par}

\medskip\noindent (5) 기타요인 \par
가. 신뢰성 있고 채택가능한 평가선례 (단위: m\textsuperscript{2}, 원/m\textsuperscript{2}) \par
{\small\centering
\begin{tabular}{|c|c|c|c|c|c|c|c|c|}
\hline
\makecell{기호} & \makecell{소재지} & \makecell{지목} & \makecell{면적} & \makecell{이용상황} & \makecell{용도\newline{}지역} & \makecell{평가단가\newline{}(가격시점)} & \makecell{평가목적} & \makecell{비고\newline{}(평가기관)} \\
\hline
\makecell{①} & \makecell{A동 159} & \makecell{대} & \makecell{250} & \makecell{상업용} & \makecell{준공업\newline{}지역} & \makecell{450,000\newline{}(2012.01.01)} & \makecell{경매} & \makecell{K법인} \\
\hline
\makecell{②} & \makecell{A동 522} & \makecell{대} & \makecell{250} & \makecell{주상용} & \makecell{준공업\newline{}지역} & \makecell{504,000\newline{}(2012.01.01)} & \makecell{보상평가} & \makecell{K법인} \\
\hline
\end{tabular}\par}

나. 비교가능하고 신뢰할 만한 최근의 거래사례는 포착하지 못하였으며, 인근 유사토지의 가격수준은 470,000원/m\textsuperscript{2}∼530,000원/m\textsuperscript{2}입니다. \par
\medskip\noindent (6) 토지가격 결정 \par
{\small\centering
\begin{tabular}{|c|c|c|c|c|c|}
\hline
\makecell{공시지가(원/m\textsuperscript{2})} & \makecell{지가변동률} & \makecell{지역요인} & \makecell{개별요인} & \makecell{기타요인} & \makecell{평가단가(원/m\textsuperscript{2})} \\
\hline
\makecell{420,000} & \makecell{1.00750} & \makecell{1.00} & \makecell{0.99} & \makecell{1.20} & \makecell{500,000} \\
\hline
\end{tabular}\par}

2) 건물가액 산출근거 \par
\medskip\noindent (1) 재조달원가: 신축건물이므로 업자가 제시한 적산자료를 활용하여 직접법에 의하여 재조달원가를 구하였습니다. \par
<표> 적산자료 (단위: 원) \par
{\small\centering
\begin{tabular}{|c|c|c|c|}
\hline
\makecell{구 분} & \makecell{내 역} & \makecell{업자제시액} & \makecell{최종사정액} \\
\hline
\makecell{설계비} & \makecell{설계비용, 감리비} & \makecell{13,000,000} & \makecell{12,000,000} \\
\hline
\makecell{기본건축비} & \makecell{기초 및 골조 공사비} & \makecell{54,000,000} & \makecell{54,000,000} \\
\hline
\makecell{옹벽공사비} & \makecell{옹벽 및 배수로 공사} & \makecell{9,000,000} & \makecell{8,000,000} \\
\hline
\makecell{내외장공사비} & \makecell{미장, 창호공사 등} & \makecell{47,000,000} & \makecell{47,000,000} \\
\hline
\makecell{위생 및 냉난방설비비} & \makecell{위생, 냉난방공사 등} & \makecell{22,000,000} & \makecell{22,000,000} \\
\hline
\makecell{전기통신설비비} & \makecell{전기 및 통신공사비 등} & \makecell{16,000,000} & \makecell{16,000,000} \\
\hline
\makecell{조경공사비} & \makecell{토지입구 조경수 등} & \makecell{14,000,000} & \makecell{14,000,000} \\
\hline
\makecell{집기 및 비품비} & \makecell{비품, 소모품 등} & \makecell{8,000,000} & \makecell{8,000,000} \\
\hline
\makecell{일반관리비 등} & \makecell{일반관리비, 이윤 등} & \makecell{15,000,000} & \makecell{14,000,000} \\
\hline
\makecell{총액} & \makecell{--} & \makecell{198,000,000} & \makecell{195,000,000} \\
\hline
\end{tabular}\par}

\medskip\noindent (2) 평가단가 \par
(단위: 년, 원/m\textsuperscript{2}) \par
{\small\centering
\begin{tabular}{|c|c|c|c|c|}
\hline
\makecell{재조달원가} & \makecell{내용년수} & \makecell{경과년수} & \makecell{산식} & \makecell{평가단가} \\
\hline
\makecell{650,000} & \makecell{50} & \makecell{0} & \makecell{650,000×50/50} & \makecell{650,000} \\
\hline
\end{tabular}\par}

4. 감정평가 명세표 \par
{\small\centering
\setlength{\tabcolsep}{3pt}
\begin{tabular}{|c|c|c|c|c|c|c|c|c|}
\hline
\multirow{2}{*}{\makecell{기호}} & \multirow{2}{*}{\makecell{소재지}} & \multirow{2}{*}{\makecell{지목\newline{}용도}} & \multirow{2}{*}{\parbox{2cm}{\centering 용도지역 및 구조}} & \multicolumn{2}{|c|}{\makecell{면적(m\textsuperscript{2})}} & \multicolumn{2}{|c|}{\makecell{평가가격}} & \multirow{2}{*}{\makecell{비고}} \\
\cline{5-8}
 &  &  &  & \makecell{공부} & \makecell{사정} & \makecell{단가} & \makecell{금액} &  \\
\hline
\makecell{①} & \makecell{K시 H구\newline{}A동 103-1} & \makecell{잡종지} & \makecell{준공업\newline{}지역} & \makecell{400} & \makecell{400} & \makecell{500,000} & \makecell{200,000,000} & \makecell{현황\newline{}“대”} \\
\hline
\makecell{㉮} & \makecell{"} & \makecell{근린생활시설} & \parbox{2cm}{\centering 철근\newline{}콘크리트조 슬래브지붕\newline{}2층} & \makecell{300} & \makecell{300} & \makecell{650,000} & \makecell{195,000,000} & \makecell{650,000× 50/50} \\
\hline
 & \makecell{합 계} &  &  &  &  &  & \makecell{395,000,000} &  \\
\hline
\end{tabular}\par}

5. 기타 부속내용(지적개황도 관련부분) \par
\qfigure{figures/02-1.png}
\textbf{<자료 2> 심사 감정평가사의 조사자료} \par
1. 각종 공적장부 확인자료 \par
토지 등기사항전부증명서를 확인한 바 103-1번지는 근저당권자 KY은행, 채권최고액은 2억원인 권리관계가 존재하고 있다. \par
2. 가격조사자료 \par
1) 경매 낙찰가율 자료(최근 6개월, 단위: \%) \par
{\small\centering
\begin{tabular}{|c|c|c|c|c|}
\hline
\makecell{지역} & \makecell{공업용} & \makecell{주거용} & \makecell{상업용} & \makecell{평균} \\
\hline
\makecell{H구} & \makecell{70} & \makecell{90} & \makecell{80} & \makecell{80} \\
\hline
\makecell{K시 전체} & \makecell{75} & \makecell{95} & \makecell{85} & \makecell{85} \\
\hline
\end{tabular}\par}

2) 평가선례 검토자료 \par
\medskip\noindent (1) 평가선례 ①: 경매 감정평가서를 열람하여 보니 평가 당시 지상의 부가물을 포함한 평균단가로 평가하였고 지상의 부가물은 조경석 3톤, 조경수 2그루였으나 구체적인 평가액은 알 수 없었다. \par
\medskip\noindent (2) 평가선례 ②: 당시 동 사업의 감정평가에 참여했던 나머지 법인의 자료가 A법인은 505,000원/m\textsuperscript{2}, B법인은 506,000원/m\textsuperscript{2}에 평가된 것으로 조사되었다. \par
3) 임대사례 조사자료 \par
\medskip\noindent (1) 공실 및 대손충당금 비율은 인근지역이 연간 가능총임대료의 15\%, 대상부동산은 5\%가 확실시 된다. \par
\medskip\noindent (2) 운영경비 비율은 모두 유사하나 시장에서 구체적인 내역을 수집하지는 못하였다. \par
\medskip\noindent (3) 대상부동산은 가격시점과 임대시점, 임대면적과 건축연면적이 각각 동일하며 연간 가능총임대료는 59,000원/m\textsuperscript{2}인 것으로 조사되었다. \par
(단위: m\textsuperscript{2}, 원) \par
{\small\centering
\begin{tabular}{|c|c|c|c|c|c|c|c|c|}
\hline
\makecell{기호} & \makecell{소재지} & \makecell{토지면적} & \makecell{임대면적} & \makecell{이용\newline{}상황} & \makecell{용도\newline{}지역} & \makecell{연간 가능총임대료} & \makecell{부동산가액} & \makecell{임대시점} \\
\hline
\makecell{①} & \makecell{A동 69} & \makecell{333} & \makecell{250} & \makecell{상업용} & \makecell{준공업\newline{}지역} & \makecell{15,000,000} & \makecell{304,000,000} & \makecell{2012.09.09} \\
\hline
\makecell{②} & \makecell{A동 90} & \makecell{400} & \makecell{300} & \makecell{상업용} & \makecell{준공업\newline{}지역} & \makecell{20,000,000} & \makecell{345,000,000} & \makecell{2012.09.09} \\
\hline
\end{tabular}\par}

\textbf{<자료 3> 기타 관련자료 등} \par
1. KY은행에서는 감정평가업무협약서상 기타요인의 산출근거를 설시하도록 요구하고 있으며, 감정평가사가 임의로 평가조건을 설정할 수 없도록 하고 있다. \par
2. 개별요인 비교자료(지역요인은 동일함) \par
{\small\centering
\begin{tabular}{|c|c|c|c|c|c|}
\hline
\makecell{대상} & \makecell{표준지 ①} & \makecell{표준지 ②} & \makecell{표준지 ③} & \makecell{평가선례 ①} & \makecell{평가선례 ②} \\
\hline
\makecell{100} & \makecell{105} & \makecell{101} & \makecell{95} & \makecell{91} & \makecell{99} \\
\hline
\end{tabular}\par}

3. 대상물건에 적용할 환원이율과 임료승수(GRM; Gross Rent Multiplier 또는 EGRM; Effective Gross Rent Multiplier)는 2개 이상의 적정한 사례자료를 기준으로 산술평균하여 구한다. \par
4. 당해 지역의 법정지상권 감안비율은 인근 송유관부지 구분지상권의 사용료(지료)가 가격에서 차지하는 비율의 10배 정도가 적정한 것으로 조사되었다. \par
1) 송유관이 지나는 105번지의 토지 등기사항전부증명서를 확인한 바 지목은 도로, 면적은 180m\textsuperscript{2}, 소유자는 乙과 丙이 각 1/2씩 소유하고 있었고 지상권자는 DH송유관공사, 사용료는 90,000원이었다. \par
2) 지상권 사용료 산정당시의 토지 감정평가액은 250,000원/m\textsuperscript{2}이었고 송유관이 지나는 부지의 폭은 3m, 길이는 6m 이었다. \par
\end{longquestion}
